%=======================================================================
%  LABORATORIO N.o 8 - TRANSISTORES DE EFECTO DE CAMPO: JFET Y MOSFET
%  Curso : Circuitos Electronicos  (2026-2)
%  UCSP  - Facultad de Ingenieria - Ing. Electronica y de Telecomunicaciones
%  Docente     : Ebert San Roman Castillo
%  Laboratorio : Ing. Jeanpier Ancori Sanchez
%
%  Version 1 - Pedido del docente: practica de JFET y MOSFET con una
%  aplicacion. Practica de DOS SEMANAS con el formato de los Laboratorios 6
%  y 7 (Datos -> Calcule -> Arme y mida -> Anote, cuadro resumen,
%  cuestionario de 3 preguntas, conclusiones 5 + 5); no se imprimen los
%  resultados.
%  Semana 1 (JFET 2N5457):
%    1. Curva de transferencia: I_DSS y V_P de la unidad de cada alumno
%       (la hoja de datos da I_DSS de 1 a 5 mA y V_GS(off) de -0,5 a -6 V).
%       Puntos en DC y la curva completa en modo XY con W1 en triangular.
%    2. Aplicacion: el JFET como fuente de corriente constante para un LED
%       (compuerta a la fuente y autopolarizado para I_DSS/2), frente a un
%       LED con solo una resistencia, con V_DD de 3, 4, 5 y 12 V.
%  Semana 2 (MOSFET 2N7000):
%    3. Umbral V_GS(th), curva de transferencia, R_DS(on) con V_GS = 5 y
%       12 V, y corriente de compuerta nula.
%    4. Aplicacion: atenuador PWM de LEDs (6 LED rojos en dos ramas en el
%       riel de 12 V). Retoma el "Lab. 08 - Circuitos accionadores de Leds"
%       del docente (2026-1), con W1 en lugar del 555 astable.
%  Instrumental: Analog Discovery Studio. Rieles fijos de +12 V (0,2 A) y
%  +5 V; fuente programable V+ para 3 y 4 V en el exp. 2. W1 da la
%  compuerta en los exp. 1, 3 y 4 (en el exp. 1 siempre negativa).
%  Consumo maximo: 36 mA en el exp. 4 y 49 mA en la medida de R_DS(on).
%
%  VALORES DE REFERENCIA PARA EL DOCENTE (a mano: I_DSS = 2,1 mA,
%  V_P = -1,4 V, V_GS(th) = 2,1 V, K = 80 mA/V^2, R_DS(on) = 2,5 Ohm con
%  4,5-5 V y 1,2 Ohm con 10-12 V, LED rojo 1,85 V a 2 mA y 2,0 V a 18 mA;
%  no se imprimen):
%    Exp. 1 (Shockley): V_GS = 0 / -0,25 / -0,5 / -0,75 / -1,0 / -1,25 V
%      -> I_D = 2,10 / 1,42 / 0,87 / 0,45 / 0,17 / 0,02 mA.
%      V_P = 2 V_GS donde I_D = I_DSS/4 (unos -0,7 V).
%    Exp. 2: I_D = I_DSS/2 = 1,05 mA ; V_GS = -1,4 (1 - 0,707) = -0,41 V ;
%      R_S = 0,41/1,05 mA = 390 Ohm (E12 exacto).
%      V_DD(min) = V_LED + |V_P| = 1,85 + 1,4 = 3,25 V (no depende de R_S).
%      Solo 4,7 k: I = (V_DD - 1,85)/4,7 k = 0,25 / 0,46 / 0,67 / 2,16 mA
%      a 3 / 4 / 5 / 12 V.
%    Exp. 3: umbral I_D = (5 - 2,1)/3,3 k = 0,88 mA.
%      Transferencia (12 V, R_D = 1 k): I_D = 80m (V_GS - 2,1)^2 hasta
%      12 mA (V_GS = 2,49 V); 1,5 y 2,0 V -> 0 ; 2,2 -> 0,8 mA ;
%      2,4 -> 7,2 mA ; 2,6 V en adelante -> 12 mA (region ohmica).
%      R_DS(on) (5 V, R_D = 100): I_D = 48,8 mA ; V_DS = 0,12 V con
%      V_GS = 5 V y 0,06 V con V_GS = 12 V. I_G = 0.
%    Exp. 4: I_on por rama = (12 - 6,0 - 0,1)/330 = 17,9 mA ; total
%      35,8 mA ; P(330) = 106 mW ; P(MOSFET) = 35,8m^2 * 2,7 = 3,5 mW.
%      I_prom = D * 35,8 mA = 3,6 / 8,9 / 17,9 / 26,8 / 32,2 mA para
%      D = 10 / 25 / 50 / 75 / 90 %. Parpadeo visible por debajo de
%      unos 50-60 Hz.
%    Cuestionario 1: V_DD(min) = V_LED + |V_P|, sin R_S (V_RS = -V_GS se
%      cancela). Con V_P = -1,4 V: 3,25 V.
%    Cuestionario 2: rele del Lab. 6 (68,6 mA) con 2N7000 a V_GS = 5 V:
%      I_G = 0 ; V_DS = 68,6m * 2,7 = 0,19 V ; P = 12,7 mW, contra
%      I_B = 3,6 mA y V_CE(sat) = 0,12 V (8,2 mW) del 2N2222A.
%    Cuestionario 3: D = 30 % -> I_prom = 10,7 mA ; P(riel) = 129 mW.
%
%  SIMULACION (ngspice 26 de Fusion/Eagle; 2N5457 de la biblioteca de
%  Motorola, 2N7000 de nivel 1 ajustado a mano a la hoja de datos y el
%  LED rojo del Lab. 6; netlists, datos y graficas en ../Simulaciones):
%    Exp. 1: I_DSS = 2,16 mA ; I_D = 2,16 / 1,45 / 0,88 / 0,45 / 0,16 /
%      0,017 mA ; I_D = 10 uA en V_GS = -1,28 V.
%    Exp. 2 (3 / 4 / 5 / 12 V): R_S = 0 -> 2,06 / 2,12 / 2,13 / 2,16 mA ;
%      R_S = 390 -> 1,03 / 1,05 / 1,05 / 1,06 mA, V_GS = -0,41 V ;
%      solo 4,7 k -> 0,28 / 0,48 / 0,69 / 2,16 mA. Con 2 V el JFET cae
%      a 0,67 mA (region ohmica).
%    Exp. 3: umbral V_GS = 2,20 V con 0,85 mA. Transferencia: 1 mA en
%      2,20 V, 10 mA en 2,45 V ; 2,2 / 2,4 / 2,6 / 3 / 5 V -> 0,92 /
%      7,72 / 11,8 / 11,9 / 12,0 mA. R_DS(on) = 2,71 Ohm (V_DS = 0,132 V,
%      48,7 mA) con 5 V y 1,21 Ohm (0,060 V, 49,4 mA) con 12 V.
%    Exp. 4: I_on = 34,6 mA (LED 2,07 V), V_DS(on) = 0,093 V ;
%      I_prom = 3,78 / 8,87 / 17,4 / 25,9 / 30,9 mA.
%=======================================================================

\documentclass[11pt,a4paper]{article}

%------------------------------------------------------------- idioma
% Compilar con XeLaTeX (no con pdflatex): Trebuchet MS es una fuente TrueType
% del sistema y solo fontspec, bajo XeLaTeX o LuaLaTeX, sabe cargarla.
\usepackage{fontspec}
\setmainfont{Trebuchet MS}[
  Path            = C:/Windows/Fonts/,
  Extension       = .ttf,
  UprightFont     = trebuc,
  BoldFont        = trebucbd,
  ItalicFont      = trebucit,
  BoldItalicFont  = trebucbi,
  Ligatures       = TeX
]
\usepackage[spanish,es-nodecimaldot]{babel}
\usepackage{microtype}

%----------------------------------------------------------- geometria
\usepackage[left=2.4cm,right=2.4cm,top=3.1cm,bottom=2.6cm,headheight=30pt,
            footskip=1.2cm]{geometry}

%------------------------------------------------------------ graficos
\usepackage{graphicx}
\usepackage{xcolor}
\usepackage{float}
\usepackage{placeins}
\usepackage[font=small,labelfont=bf,labelsep=period]{caption}
\usepackage{subcaption}
\graphicspath{{imagenes/}}
\usepackage[american,siunitx]{circuitikz}
% babel-spanish activa < > " y rompe las opciones de TikZ sin esta libreria.
\usetikzlibrary{calc,patterns,babel,arrows.meta,positioning,decorations.pathreplacing}

%-------------------------------------------------------------- tablas
\usepackage{array}
\usepackage{tabularx}
\usepackage{multirow}

%----------------------------------------------------- matematicas / SI
\usepackage{amsmath,amssymb}
\usepackage{siunitx}
\sisetup{
  output-decimal-marker = {,},
  per-mode              = symbol,
  mode                  = match,
  detect-weight         = true,
  detect-family         = true
}

%--------------------------------------------------------------- listas
\usepackage{enumitem}
\setlist[itemize]{itemsep=3pt,topsep=4pt,parsep=0pt}
\setlist[enumerate]{itemsep=8pt,topsep=4pt,parsep=0pt}

%------------------------------------------------------------ comandos
\usepackage{pgffor}

% Lineas punteadas para respuestas del alumno: \lineasresp{n}
\newcommand{\lineasresp}[1]{%
  \par\vspace{6pt}%
  \foreach \i in {1,...,#1}{\noindent\rule{\linewidth}{0.4pt}\par\vspace{8pt}}%
}

% Recuadro "Figura / Comentario": una celda alta para pegar la captura del
% osciloscopio y otra debajo para el comentario.
\newcommand{\cuadrorespuesta}{%
  \par\vspace{6pt}\noindent
  \begin{tabularx}{\linewidth}{|>{\centering\arraybackslash}X|}
    \hline
    \rule{0pt}{4.4cm}\raisebox{2.1cm}{\textcolor{black!45}{Figura}}\\
    \hline
    \rule{0pt}{2.0cm}\raisebox{0.9cm}{\textcolor{black!45}{Comentario}}\\
    \hline
  \end{tabularx}\par\vspace{8pt}%
}

% Celda vacia de altura comoda para escribir a mano.
\newcommand{\hueco}{\rule{0pt}{0.5cm}}

%----------------------------------------------------- compaginacion
\usepackage{needspace}
% \needspace mide mal el espacio libre despues de figuras y tablas [H] y deja
% paginas medio vacias; \Needspace es la variante exacta del mismo paquete.
\let\needspace\Needspace
\usepackage{etoolbox}
\pretocmd{\section}{\needspace{5\baselineskip}}{}{}
\pretocmd{\subsection}{\needspace{4\baselineskip}}{}{}
\widowpenalty=10000
\clubpenalty=10000

%---------------------------------------------------- encabezado / pie
\usepackage{fancyhdr}
\pagestyle{fancy}
\fancyhf{}
\fancyhead[L]{\footnotesize UCSP --- Facultad de Ingeniería\\
              Ing. Electrónica y de Telecomunicaciones}
\fancyhead[R]{\footnotesize 2026-2 Circuitos Electrónicos\\
              Ebert San Román Castillo}
\fancyfoot[L]{\footnotesize Laboratorio N.\textsuperscript{o} 8 --- Transistores de efecto de campo: JFET y MOSFET}
\fancyfoot[R]{\footnotesize Página \thepage\ de \pageref{LastPage}}
\renewcommand{\headrulewidth}{0.4pt}
\renewcommand{\footrulewidth}{0.4pt}
\usepackage{lastpage}

%------------------------------------------------------- numeracion
\renewcommand{\thesection}{\Roman{section}}
\renewcommand{\thesubsection}{\arabic{subsection}}
\renewcommand{\thesubsubsection}{\thesubsection.\arabic{subsubsection}}
\setcounter{secnumdepth}{3}

%-------------------------------------------------------- hiperenlaces
\usepackage[hidelinks,pdfusetitle]{hyperref}
\usepackage{xurl}
\hypersetup{
  pdftitle={Laboratorio 8 - Transistores de efecto de campo: JFET y MOSFET},
  pdfsubject={Circuitos Electronicos - Guia de practicas},
  pdfauthor={Ebert San Roman Castillo; Jeanpier Ancori Sanchez},
  pdfkeywords={electronica, transistor de efecto de campo, FET, JFET, MOSFET,
               2N5457, 2N7000, curva de transferencia, ecuacion de Shockley,
               IDSS, voltaje de estrangulamiento, voltaje umbral, RDS(on),
               autopolarizacion, fuente de corriente, PWM, LED,
               Analog Discovery Studio, Proteus, Multisim, osciloscopio}
}

% babel-spanish rotula las tablas como "Cuadro"; los originales dicen "Tabla".
\addto\captionsspanish{\renewcommand{\tablename}{Tabla}}

\begin{document}

%=======================================================================
%  CABECERA DE LA GUIA
%=======================================================================
\begin{center}
    {\normalsize\bfseries SEGUNDA UNIDAD: TRANSISTOR DE EFECTO DE CAMPO (FET)}
\end{center}

\vspace{0.55cm}

\begingroup
\renewcommand{\tabularxcolumn}[1]{m{#1}}
\renewcommand{\arraystretch}{2.0}
\setlength{\tabcolsep}{12pt}
\noindent
\begin{tabularx}{\textwidth}{|>{\centering\arraybackslash}m{2.2cm}|>{\centering\arraybackslash}X|>{\centering\arraybackslash}m{3.0cm}|}
\hline
\rule{0pt}{2.1cm}{\fontsize{58}{58}\selectfont\bfseries 8}
    & \begin{minipage}[c]{\linewidth}\centering
          {\large\bfseries Guía de Prácticas}\\[7pt]
          {\Large\bfseries Transistores de Efecto\\ de Campo: JFET y MOSFET}
      \end{minipage}
    & {\Large\bfseries Nota}\\
\hline
\multicolumn{3}{|l|}{\rule{0pt}{1.3cm}\large\textbf{Grupo:}~\rule{4.6cm}{0.4pt}%
    \hfill\large\textbf{Fecha:}~\rule{4.6cm}{0.4pt}\hspace{2mm}}\\
\hline
\multicolumn{3}{|l|}{\rule{0pt}{1.2cm}\large\textbf{Alumno(s):}~\rule{11.2cm}{0.4pt}\hspace{2mm}}\\
\multicolumn{3}{|l|}{\rule{0pt}{1.2cm}\phantom{\large\textbf{Alumno(s):}}~\rule{11.2cm}{0.4pt}\hspace{2mm}}\\
\hline
\end{tabularx}
\endgroup

\vspace{0.8cm}

% Encabezado de cada paso dentro de un experimento: \paso{1}{Datos}
\newcommand{\paso}[2]{%
  \par\needspace{6\baselineskip}\vspace{8pt}\noindent
  \colorbox{black!8}{\makebox[\dimexpr\linewidth-2\fboxsep\relax][l]{%
    \textbf{Paso #1 \,·\, #2}}}\par\vspace{4pt}}

%=======================================================================
\section{Objetivos}
%=======================================================================
\begin{itemize}
    \item Medir la curva de transferencia de un JFET, obtener su corriente $I_{DSS}$ y su
          voltaje de estrangulamiento $V_P$, y compararlos con la ecuación de Shockley.
    \item Diseñar la autopolarización de un JFET y usarlo como fuente de corriente constante
          para un LED, comprobando que la corriente no cambia con el voltaje de alimentación.
    \item Medir el voltaje umbral, la curva de transferencia y la resistencia de conducción
          de un MOSFET de enriquecimiento, y usarlo como interruptor en un atenuador PWM de
          LEDs.
\end{itemize}

%=======================================================================
\section{Contenido teórico}
%=======================================================================
\begin{itemize}
    \item El transistor de efecto de campo (FET): control por voltaje, corriente de compuerta
          nula; JFET y MOSFET.
    \item El JFET: ecuación de Shockley, regiones de trabajo, autopolarización y fuente de
          corriente constante.
    \item El MOSFET de enriquecimiento: voltaje umbral, región de saturación y región óhmica;
          el MOSFET como interruptor y la modulación por ancho de pulso (PWM).
\end{itemize}

\noindent\textbf{Duración:} dos sesiones. \textbf{Semana~1:} experimentos~1 y~2 (JFET).
\textbf{Semana~2:} experimentos~3 y~4 (MOSFET).

%=======================================================================
\section{Equipos y materiales}
%=======================================================================
\noindent\textbf{Laboratorio:} Electrónica de Comunicaciones.

\vspace{0.25cm}
\noindent\textbf{Equipos y dispositivos:}
\begin{itemize}
    \item Tarjeta Analog Discovery Studio: rieles fijos de \qty{+12}{\volt} y
          \qty{+5}{\volt}, fuente programable $V_+$, generador de funciones \texttt{W1} y
          osciloscopio de dos canales.
\end{itemize}

\noindent\textbf{Software:}
\begin{itemize}
    \item Simulador Multisim o Proteus.
\end{itemize}

\vspace{0.25cm}
\needspace{6\baselineskip}
\noindent\textbf{Materiales y fungibles} (deben ser adquiridos por los alumnos):
\begin{itemize}
    % --- Semiconductores
    \item 02 JFET canal N 2N5457 (o MPF102) y 02 MOSFET canal N 2N7000 (uno de repuesto de
          cada uno). Sus terminales están en distinto orden: revise la hoja de datos antes de
          armar.
    \item 08 LED rojos de \qty{5}{\milli\metre} (uno de repuesto).
    % --- Resistencias (la que mas disipa es RD = 100 Ohm del exp. 3: 240 mW;
    %     las de 330 Ohm del exp. 4 reciben 106 mW)
    \item Resistencia a \qty{0.5}{\watt}: 01 de \qty{100}{\ohm}.
    \item Resistencias a \qty{0.25}{\watt}:
          01 de \qty{470}{\ohm},
          02 de \qty{100}{\kilo\ohm},
          01 de \qty{1}{\mega\ohm},
          01 de \qty{4.7}{\kilo\ohm},
          01 de \qty{3.3}{\kilo\ohm},
          02 de \qty{1}{\kilo\ohm} y
          02 de \qty{330}{\ohm}.
          Para la $R_S$ del experimento~2 traiga una de cada valor:
          \qty{220}{\ohm}, \qty{270}{\ohm}, \qty{390}{\ohm}, \qty{560}{\ohm},
          \qty{680}{\ohm} y \qty{820}{\ohm}.
    % --- Montaje
    \item Protoboard.
    \item Jumpers macho-macho para protoboard.
    \item Alicates para cortar y pelar alambre.
    \item Multímetro.
    \item Lentes de protección y mandil.
\end{itemize}

\noindent La compuerta del MOSFET es una capa de óxido muy delgada que se perfora con la
electricidad estática: tómelo por la cápsula y no deje nunca su compuerta al aire.

\FloatBarrier
%=======================================================================
\section{Conocimientos previos}
%=======================================================================

%-----------------------------------------------------------------------
\subsection{El transistor de efecto de campo}
%-----------------------------------------------------------------------
En el transistor bipolar (Laboratorios~5 a~7) una \textbf{corriente} de base controla la
corriente de colector. En el transistor de efecto de campo (FET) un \textbf{voltaje}, el
voltaje compuerta-fuente $V_{GS}$, controla la corriente que circula por un canal entre el
\textbf{drenador} (D) y la \textbf{fuente} (S). La \textbf{compuerta} (G) está aislada del
canal, así que no toma corriente: $I_G \approx 0$ y la resistencia de entrada es de
\qty{e9}{\ohm} o más. Hay dos familias (figura~\ref{fig:simbolos}):
\begin{itemize}
  \item \textbf{JFET} (FET de unión): la compuerta es una unión PN polarizada en
        \emph{inversa}. Conduce con $V_{GS} = 0$ y se va cerrando con $V_{GS}$ negativo
        (canal N). Es un dispositivo de \textbf{agotamiento}.
  \item \textbf{MOSFET} (FET de compuerta aislada): la compuerta está separada del canal por
        una capa de óxido. El de \textbf{enriquecimiento}, el más usado, no conduce con
        $V_{GS} = 0$: el canal aparece cuando $V_{GS}$ supera un umbral positivo (canal N).
\end{itemize}

\begin{figure}[H]
  \centering
  \begin{circuitikz}[scale=0.9, transform shape]
    \draw (0,0) node[npn](Q){};
    \node[right] at (Q.C) {C};  \node[left] at (Q.B) {B};  \node[right] at (Q.E) {E};
    \node[below=6mm, font=\small] at (Q.E) {BJT NPN};
    \draw (4,0) node[njfet](J){};
    \node[right] at (J.D) {D};  \node[left] at (J.G) {G};  \node[right] at (J.S) {S};
    \node[below=6mm, font=\small] at (J.S) {JFET canal N};
    \draw (8,0) node[nigfete](M){};
    \node[right] at (M.D) {D};  \node[left] at (M.G) {G};  \node[right] at (M.S) {S};
    \node[below=6mm, font=\small] at (M.S) {MOSFET canal N (enriq.)};
  \end{circuitikz}
  \caption{Símbolos. El canal partido del MOSFET indica que no existe hasta que $V_{GS}$
           supera el umbral.}
  \label{fig:simbolos}
\end{figure}

%-----------------------------------------------------------------------
\subsection{El JFET: ecuación de Shockley}
\label{sec:jfet}
%-----------------------------------------------------------------------
Con $V_{DS}$ suficientemente grande el JFET trabaja en la \textbf{región de saturación}
(o de estrangulamiento): la corriente de drenador ya no depende de $V_{DS}$, solo de $V_{GS}$,
y sigue la \textbf{ecuación de Shockley} (figura~\ref{fig:transf}a):
\begin{equation}
  I_D = I_{DSS}\left(1 - \frac{V_{GS}}{V_P}\right)^{2},
  \qquad V_P \le V_{GS} \le 0.
  \label{eq:shockley}
\end{equation}
$I_{DSS}$ es la corriente con la compuerta unida a la fuente ($V_{GS} = 0$), la máxima del
JFET. $V_P$ (o $V_{GS(\text{off})}$) es el voltaje de \textbf{estrangulamiento}: con
$V_{GS} \le V_P$ el canal se cierra y $I_D = 0$ (\textbf{corte}). Las tres regiones son:
\begin{itemize}
  \item \textbf{Corte:} $V_{GS} \le V_P$; no hay corriente.
  \item \textbf{Saturación:} $V_{DS} \ge V_{GS} - V_P$; $I_D$ sigue~\eqref{eq:shockley}.
        Es la región de los amplificadores y de la fuente de corriente.
  \item \textbf{Óhmica:} $V_{DS} < V_{GS} - V_P$; el canal se comporta como una resistencia
        controlada por $V_{GS}$.
\end{itemize}

\noindent\textbf{Cómo medir $V_P$.} Cerca de $V_P$ la corriente es tan pequeña que cuesta
ver dónde se anula. Es más preciso usar el punto en que $I_D = I_{DSS}/4$: por~\eqref{eq:shockley}
ahí $1 - V_{GS}/V_P = 1/2$, es decir,
\begin{equation}
  V_P = 2\,V_{GS} \quad\text{cuando}\quad I_D = \frac{I_{DSS}}{4}.
  \label{eq:vp}
\end{equation}

\noindent\textbf{Dispersión.} $I_{DSS}$ y $V_P$ cambian mucho de una unidad a otra: la hoja de
datos del 2N5457 da $I_{DSS}$ de \qtyrange{1}{5}{\milli\ampere} y $V_{GS(\text{off})}$ de
\qtyrange{-0.5}{-6}{\volt}. Por eso cada alumno mide los de su propio JFET (experimento~1)
antes de diseñar con él (experimento~2).

\begin{figure}[H]
  \centering
  \begin{subfigure}[b]{0.48\linewidth}
    \centering
    \begin{tikzpicture}[scale=0.95, >=Stealth]
      \draw[->] (-3.6,0) -- (0.6,0) node[below] {$V_{GS}$};
      \draw[->] (0,-0.2) -- (0,3.4) node[left] {$I_D$};
      \draw[very thick, domain=-3:0, smooth, samples=40, variable=\x]
           plot ({\x},{2.8*(1+\x/3)^2});
      \draw[dashed] (-1.5,0) node[below] {$V_P/2$} -- (-1.5,0.7) -- (0,0.7)
           node[right] {$I_{DSS}/4$};
      \fill (-1.5,0.7) circle (1.6pt);
      \fill (0,2.8) circle (1.6pt) node[right] {$I_{DSS}$};
      \node[below] at (-3,0) {$V_P$};
      \fill (-3,0) circle (1.6pt);
      \node[font=\small, align=center] at (-2.6,2.4) {corte a la\\izquierda de $V_P$};
    \end{tikzpicture}
    \caption{JFET canal N.}
  \end{subfigure}
  \hfill
  \begin{subfigure}[b]{0.48\linewidth}
    \centering
    \begin{tikzpicture}[scale=0.95, >=Stealth]
      \draw[->] (-0.2,0) -- (4.2,0) node[below] {$V_{GS}$};
      \draw[->] (0,-0.2) -- (0,3.4) node[left] {$I_D$};
      \draw[very thick] (0,0) -- (1.4,0);
      \draw[very thick, domain=1.4:2.6, smooth, samples=30, variable=\x]
           plot ({\x},{2.0*(\x-1.4)^2});
      \draw[very thick] (2.6,2.88) .. controls (2.75,3.0) .. (3.9,3.05);
      \fill (1.4,0) circle (1.6pt);
      \node[below] at (1.4,0) {$V_{GS(\text{th})}$};
      \node[font=\small, align=center] at (3.2,1.6) {la carga\\limita $I_D$};
      \node[font=\small] at (0.7,0.45) {corte};
    \end{tikzpicture}
    \caption{MOSFET canal N de enriquecimiento.}
  \end{subfigure}
  \caption{Curvas de transferencia $I_D$ contra $V_{GS}$. El JFET conduce con $V_{GS} = 0$ y se
           corta con $V_{GS}$ negativo; el MOSFET de enriquecimiento no conduce hasta que
           $V_{GS}$ supera el umbral.}
  \label{fig:transf}
\end{figure}

%-----------------------------------------------------------------------
\subsection{Autopolarización del JFET}
%-----------------------------------------------------------------------
Para trabajar con un $V_{GS}$ negativo sin una fuente negativa se pone una resistencia $R_S$
en la fuente y la compuerta a tierra a través de $R_G$ (figura~\ref{fig:exp2}b). Como
$I_G = 0$, la compuerta queda a \qty{0}{\volt} y la fuente a $I_D R_S$, así que
\begin{equation}
  V_{GS} = -I_D\,R_S.
  \label{eq:autopol}
\end{equation}
\noindent\textbf{Cómo se diseña.} Para una corriente deseada $I_D$, se despeja $V_{GS}$
de~\eqref{eq:shockley} y luego $R_S$ de~\eqref{eq:autopol}:
\begin{equation}
  V_{GS} = V_P\left(1 - \sqrt{\frac{I_D}{I_{DSS}}}\right),
  \qquad
  R_S = \frac{-V_{GS}}{I_D}.
  \label{eq:rs}
\end{equation}
Con el valor E12 más cercano se recalcula $I_D$ y se comprueba que el JFET esté en saturación:
$V_{DS} \ge V_{GS} - V_P$.

\needspace{14\baselineskip}
\noindent\textbf{Ejemplo resuelto} (valores distintos a los de la práctica). JFET con
$I_{DSS} = \qty{8}{\milli\ampere}$ y $V_P = \qty{-4}{\volt}$, $V_{DD} = \qty{15}{\volt}$,
$R_D = \qty{2.2}{\kilo\ohm}$; se desea $I_D = I_{DSS}/4 = \qty{2}{\milli\ampere}$.
\begin{table}[H]
  \centering
  \renewcommand{\arraystretch}{2.1}
  \begin{tabularx}{\linewidth}{|>{\centering\arraybackslash}m{2.6cm}|>{\centering\arraybackslash}X|>{\centering\arraybackslash}m{2.8cm}|}
    \hline
    \textbf{Paso} & \textbf{Cálculo} & \textbf{Resultado} \\ \hline
    $V_{GS}$ & $-4\left(1 - \sqrt{\dfrac{2}{8}}\right) = -4\,(1 - \num{0.5})$
             & \qty{-2}{\volt} \\ \hline
    $R_S$ & $\dfrac{\qty{2}{\volt}}{\qty{2}{\milli\ampere}}$ & \qty{1}{\kilo\ohm} (E12) \\ \hline
    $V_{DS}$ & $15 - \qty{2}{\milli\ampere}\,(\num{2.2} + 1)\,\unit{\kilo\ohm}$
             & \qty{8.6}{\volt} \\ \hline
    Comprobación & $V_{GS} - V_P = -2 - (-4) = \qty{2}{\volt}$, \qquad $\num{8.6} \ge 2$
                 & saturación \\ \hline
  \end{tabularx}
  \caption{Ejemplo de autopolarización de un JFET.}
  \label{tab:ejemplo}
\end{table}

%-----------------------------------------------------------------------
\subsection{El JFET como fuente de corriente constante}
\label{sec:fuente}
%-----------------------------------------------------------------------
En la región de saturación $I_D$ depende de $V_{GS}$ pero \textbf{no} de $V_{DS}$. Si se pone
la carga (un LED) en el drenador, la corriente queda fijada por el JFET y su $R_S$, aunque cambie
el voltaje de alimentación: el JFET absorbe la diferencia en su $V_{DS}$. Con la compuerta
unida a la fuente la corriente es $I_{DSS}$; con una $R_S$, una corriente menor dada
por~\eqref{eq:rs}. Esto funciona mientras el JFET siga en saturación; el voltaje mínimo de
alimentación es
\begin{equation}
  V_{DD(\min)} = V_{\text{LED}} + I_D R_S + (V_{GS} - V_P).
  \label{eq:vddmin}
\end{equation}
En cambio, con un LED y solo una resistencia $R$, la corriente
$I = (V_{DD} - V_{\text{LED}})/R$ sigue cada variación de $V_{DD}$.

%-----------------------------------------------------------------------
\subsection{El MOSFET de enriquecimiento}
%-----------------------------------------------------------------------
Con $V_{GS} = 0$ no hay canal y el MOSFET está en \textbf{corte}. Cuando $V_{GS}$ supera el
\textbf{voltaje umbral} $V_{GS(\text{th})}$ (entre \num{0.8} y \qty{3}{\volt} en el 2N7000) el
canal se forma y, en la región de saturación,
\begin{equation}
  I_D = K\,\bigl(V_{GS} - V_{GS(\text{th})}\bigr)^{2},
  \qquad V_{GS} > V_{GS(\text{th})},
  \label{eq:mos}
\end{equation}
donde $K$ (en \unit{\milli\ampere\per\volt\squared}) depende del transistor. La hoja de datos
define $V_{GS(\text{th})}$ como el $V_{GS}$ con el que circula \qty{1}{\milli\ampere} cuando la
compuerta está unida al drenador ($V_{DS} = V_{GS}$).

\noindent\textbf{El MOSFET como interruptor.} Con $V_{GS}$ muy por encima del umbral el
MOSFET pasa a la \textbf{región óhmica} y se comporta como una resistencia pequeña, la
\textbf{resistencia de conducción} $R_{DS(\text{on})}$:
\begin{equation}
  V_{DS} = I_D\,R_{DS(\text{on})}.
  \label{eq:rdson}
\end{equation}
$R_{DS(\text{on})}$ baja cuanto mayor es $V_{GS}$. La tabla~\ref{tab:bjtmos} compara este
interruptor con el del Laboratorio~6.

\begin{table}[H]
  \centering
  \renewcommand{\arraystretch}{1.3}
  \begin{tabularx}{\linewidth}{|>{\raggedright\arraybackslash}m{3.6cm}|>{\centering\arraybackslash}X|>{\centering\arraybackslash}X|}
    \hline
    & \textbf{BJT (Laboratorio 6)} & \textbf{MOSFET} \\ \hline
    Se controla con & corriente de base $I_B$ & voltaje $V_{GS}$ \\ \hline
    Corriente de entrada en DC & $I_C/\beta_{\text{forzado}}$ & $\approx 0$ \\ \hline
    Caída en conducción & $V_{CE(\text{sat})} \approx \qty{0.2}{\volt}$ & $I_D\,R_{DS(\text{on})}$ \\ \hline
    Para conducir necesita & $V_{BE} \approx \qty{0.7}{\volt}$ e $I_B$ suficiente & $V_{GS}$ muy por encima de $V_{GS(\text{th})}$ \\ \hline
    Cuidado & $\beta$ distinto en cada unidad & compuerta sensible a la estática \\ \hline
  \end{tabularx}
  \caption{El transistor bipolar y el MOSFET como interruptores.}
  \label{tab:bjtmos}
\end{table}

%-----------------------------------------------------------------------
\subsection{Modulación por ancho de pulso (PWM)}
%-----------------------------------------------------------------------
Para variar el brillo de un LED sin desperdiciar potencia en una resistencia variable, se lo
enciende y apaga muy rápido con un interruptor. La señal de control es una onda cuadrada de
período $T$ que está en alto un tiempo $t_{on}$ (figura~\ref{fig:pwm}). Su \textbf{ciclo de
trabajo} es $D = t_{on}/T$ y la corriente promedio en la carga es
\begin{equation}
  I_{\text{prom}} = D\,I_{on},
  \label{eq:pwm}
\end{equation}
\noindent\begin{minipage}[c]{0.50\linewidth}
donde $I_{on}$ es la corriente con el interruptor cerrado. Por encima de unos
\qty{100}{\hertz} el ojo no distingue el parpadeo y percibe un brillo que crece con $D$. El
MOSFET disipa poco: casi nada cuando está abierto y solo $I_{on}^2 R_{DS(\text{on})}$ cuando
está cerrado.
\end{minipage}\hfill
\begin{minipage}[c]{0.46\linewidth}
  \centering
  \begin{tikzpicture}[scale=0.62, >=Stealth]
    \draw[->] (0,0) -- (7.4,0) node[below] {$t$};
    \draw[->] (0,0) -- (0,2.0) node[left] {$v_{GS}$};
    \draw[very thick] (0,0) -- (0,1.3) -- (1.2,1.3) -- (1.2,0) -- (3.2,0) -- (3.2,1.3)
          -- (4.4,1.3) -- (4.4,0) -- (6.4,0) -- (6.4,1.3) -- (7.0,1.3);
    \node[left] at (0,1.3) {\qty{5}{\volt}};
    \draw[<->] (3.2,1.65) -- node[above] {$t_{on}$} (4.4,1.65);
    \draw[<->] (3.2,-0.4) -- node[below] {$T$} (6.4,-0.4);
  \end{tikzpicture}
  \captionof{figure}{Señal PWM con $D = \qty{37.5}{\percent}$. El LED recibe $I_{on}$ durante
                     $t_{on}$ y nada el resto del período.}
  \label{fig:pwm}
\end{minipage}

\FloatBarrier
\needspace{14\baselineskip}
%=======================================================================
\section{Desarrollo del laboratorio}
%=======================================================================
La práctica dura dos semanas. Cada experimento tiene los mismos cuatro pasos. Los pasos~1 y~2
(\textbf{Datos} y \textbf{Calcule}) se traen resueltos de casa, junto con la simulación en
Multisim o Proteus; los pasos~3 y~4 (\textbf{Arme y mida} y \textbf{Anote}) se hacen en el
laboratorio.

\begin{table}[H]
  \centering
  \small
  \renewcommand{\arraystretch}{1.2}
  \begin{tabularx}{\linewidth}{|>{\centering\arraybackslash}m{0.9cm}|>{\centering\arraybackslash}m{0.8cm}|>{\raggedright\arraybackslash}m{3.0cm}|>{\raggedright\arraybackslash}X|>{\raggedright\arraybackslash}X|>{\centering\arraybackslash}m{1.3cm}|}
    \hline
    \textbf{Sem.} & \textbf{Exp.} & \textbf{Circuito} & \textbf{En casa: calcule}
      & \textbf{En el laboratorio: mida} & \textbf{Fuente} \\ \hline
    \multirow{2}{*}{1}
      & 1 & Curva del JFET
          & $I_D$ con la ecuación de Shockley
          & $I_D$ contra $V_{GS}$; $I_{DSS}$ y $V_P$; curva en modo XY
          & \qty{12}{\volt} \\ \cline{2-6}
      & 2 & JFET como fuente de corriente
          & $R_S$, $V_{GS}$ y $V_{DD(\min)}$; $I$ con solo $R$
          & $I_D$ con \qty{3}{} a \qty{12}{\volt}; punto Q
          & $V_+$ y \qty{12}{\volt} \\ \hline
    \multirow{2}{*}{2}
      & 3 & MOSFET 2N7000
          & $I_D$ en el umbral y contra $V_{GS}$; $V_{DS}$ en conducción
          & $V_{GS(\text{th})}$; curva en modo XY; $R_{DS(\text{on})}$
          & \qty{12}{} y \qty{5}{\volt} \\ \cline{2-6}
      & 4 & Atenuador PWM de LEDs
          & $I_{on}$, potencias e $I_{\text{prom}}$
          & $I_{\text{prom}}$ y brillo contra $D$; parpadeo
          & \qty{12}{\volt} \\ \hline
  \end{tabularx}
  \caption{Resumen de la práctica.}
  \label{tab:resumen}
\end{table}

%-----------------------------------------------------------------------
\subsection{Experimento 1: curva de transferencia del JFET}
%-----------------------------------------------------------------------
Se mide cómo $V_{GS}$ controla la corriente de drenador y se obtienen el $I_{DSS}$ y el $V_P$
del JFET de cada alumno.

\noindent\begin{minipage}{\linewidth}
\paso{1}{Datos}
\noindent\begin{minipage}[c]{0.46\linewidth}
  \centering
  \renewcommand{\arraystretch}{1.2}
  \begin{tabularx}{\linewidth}{|>{\raggedright\arraybackslash}X|>{\centering\arraybackslash}m{1.9cm}|}
    \hline
    $V_{DD}$ & \qty{12}{\volt} \\ \hline
    $R_D$ & \qty{470}{\ohm} \\ \hline
    $R_G$ & \qty{100}{\kilo\ohm} \\ \hline
    JFET & 2N5457 \\ \hline
    $I_{DSS}$ (en casa) & \qty{2.1}{\milli\ampere} \\ \hline
    $V_P$ (en casa) & \qty{-1.4}{\volt} \\ \hline
  \end{tabularx}
  \captionof{table}{Experimento 1: datos. $I_{DSS}$ y $V_P$ son los del modelo de simulación;
                    los de su JFET los medirá aquí.}
  \label{tab:cond1}
\end{minipage}\hfill
\begin{minipage}[c]{0.50\linewidth}
  \centering
  \begin{circuitikz}[scale=0.85, transform shape]
    \draw (3.2,0) node[njfet](J){};
    \node[right] at ($(J.S)+(0.05,0.35)$) {2N5457};
    \draw (J.D) -- ++(0,0.3) coordinate(d) to[R, l=$R_D$] ++(0,2.0) coordinate(top);
    \draw (top) node[circ]{} -- ++(0.8,0) node[right]{$V_{DD}$};
    \draw (d) to[short, *-o] ++(1.0,0) node[right]{$D$};
    \draw (J.G) to[R, l_=$R_G$] ++(-2.0,0) node[ocirc]{} node[left]{\texttt{W1}};
    \draw (J.S) node[ground]{};
  \end{circuitikz}
  \captionof{figure}{Experimento 1. Como $I_G \approx 0$, en $R_G$ no cae voltaje y $V_{GS}$
                     es el voltaje de \texttt{W1}.}
  \label{fig:exp1}
\end{minipage}
\end{minipage}\par\vspace{6pt}

\paso{2}{Calcule (en casa)}
\begin{enumerate}[label=\alph*), leftmargin=*, itemsep=3pt]
  \item $I_D$ con~\eqref{eq:shockley} para los $V_{GS}$ de la tabla~\ref{tab:exp1}, y el
        voltaje en $R_D$ que corresponde a cada uno.
  \item Compruebe que con $V_{GS} = 0$ el JFET está en saturación:
        $V_{DS} \ge V_{GS} - V_P$.
  \item Simule el circuito con un barrido DC de $V_{GS}$ de \qty{0}{\volt} a
        \qty{-1.6}{\volt}.
\end{enumerate}

% Conexion con la tarjeta, en los colores del arnes de la Analog Discovery
% Studio (W1 amarillo, osciloscopio 1 naranja y 2 azul, GND negro,
% alimentacion roja). Los canales son diferenciales: CH2 va sobre R_D.
\begin{figure}[H]
  \centering
  \begin{circuitikz}[scale=0.85, transform shape]
    \colorlet{cw1}{yellow!70!black}
    \colorlet{cch1}{orange!90!black}
    \colorlet{cch2}{blue!75!black}
    \colorlet{cvcc}{red!80!black}
    \coordinate (gnd) at (0,-1.6);
    \draw (4.2,0) node[njfet](J){};
    \node[right] at ($(J.S)+(0.05,0.35)$) {2N5457};
    \draw (J.D) -- ++(0,0.3) coordinate(d) to[R, l_=$R_D$] ++(0,2.0) coordinate(top);
    \draw (top) -- ++(0,0.6) coordinate(vtop);
    \coordinate (a) at ($(J.G)+(-2.4,0)$);
    \draw (a) to[R, l=$R_G$, *-] (J.G);
    \draw (J.S) -- (J.S |- gnd) node[circ]{};
    % bus de tierra
    \coordinate (p2m) at (6.2,0);
    \coordinate (p2p) at (7.0,0);
    \coordinate (p12) at (8.6,0);
    \draw[very thick] (-0.8,-1.6) -- (J.S |- gnd);
    % tarjeta
    \coordinate (ytab) at (0,-3.2);
    \draw[thick, fill=black!6, rounded corners=3pt]
         (-1.6,-3.2) rectangle (9.4,-4.7);
    \node[font=\bfseries] at (3.9,-4.25) {Analog Discovery Studio};
    \coordinate (pw1) at ($(a |- ytab)+(-1.0,0)$);
    \coordinate (p1p) at (a |- ytab);
    \coordinate (p1m) at ($(a |- ytab)+(0.6,0)$);
    \fill[black] (-0.92,-3.32) rectangle (-0.68,-3.08);
    \node[font=\scriptsize\ttfamily, below] at (-0.8,-3.32) {GND};
    \foreach \p/\n/\c in {pw1/W1/cw1, p1p/1+/cch1, p1m/1$-$/cch1,
                          p2m/2$-$/cch2, p2p/2+/cch2, p12/+12V/cvcc}{
      \fill[\c] ($(\p |- ytab)+(-0.12,-0.12)$) rectangle ($(\p |- ytab)+(0.12,0.12)$);
      \node[font=\scriptsize\ttfamily, below] at ($(\p |- ytab)+(0,-0.12)$) {\n};
    }
    % cables
    \draw[black, line width=1.3pt] (-0.8,-3.08) -- (-0.8,-1.6) node[circ]{};
    \draw[cw1, line width=1.3pt]  ($(pw1)+(0,0.12)$) -- (pw1 |- a) -- (a);
    \draw[cch1, line width=1.3pt] ($(p1p)+(0,0.12)$) -- (a);
    \draw[cch1, line width=1.3pt, dashed] ($(p1m)+(0,0.12)$) -- (p1m |- gnd) node[circ, black]{};
    \draw[cch2, line width=1.3pt] ($(p2m |- ytab)+(0,0.12)$) -- (p2m |- d) -- (d) node[circ, black]{};
    \draw[cch2, line width=1.3pt] ($(p2p |- ytab)+(0,0.12)$) -- (p2p |- top) -- (top) node[circ, black]{};
    \draw[cvcc, line width=1.3pt] ($(p12 |- ytab)+(0,0.12)$) -- (p12 |- vtop) -- (vtop);
    % rotulos
    \node[cch1, font=\scriptsize, above] at ($(a)+(-0.5,0)$) {CH1: $v_{GS}$};
    \node[cch2, font=\scriptsize, right] at ($(p2p |- top)!0.5!(p2p |- d)$) {CH2: $v_{R_D}$};
  \end{circuitikz}
  \caption{Conexión con la Analog Discovery Studio. \texttt{W1} da $V_{GS}$ y \texttt{1+} lo
           mide. CH2 mide sobre $R_D$: \texttt{2+} arriba y \texttt{2$-$} en el drenador, así
           que $I_D = v_{\text{CH2}}/R_D$. El experimento~3 usa la misma conexión con el
           MOSFET y $R_D = \qty{1}{\kilo\ohm}$.}
  \label{fig:conexion}
\end{figure}

\paso{3}{Arme y mida (en el laboratorio)}
\begin{enumerate}[label=\alph*), ref=\alph*, leftmargin=*, itemsep=3pt]
  \item Arme el circuito y conecte la tarjeta como en la figura~\ref{fig:conexion}.
        \textbf{El voltaje de \texttt{W1} debe ser siempre negativo o cero}: con la
        compuerta positiva la unión compuerta-canal conduce.
  \item Ponga \texttt{W1} en continua con cada $V_{GS}$ de la tabla~\ref{tab:exp1} y mida el
        voltaje en $R_D$ (CH2); $I_D = v_{R_D}/\qty{470}{\ohm}$. Si con
        \qty{-1.25}{\volt} todavía circula corriente, siga con \qty{-1.5}{\volt},
        \qty{-2}{\volt}, \qty{-3}{\volt}\dots{} hasta que $I_D$ sea casi cero.
  \item Con $I_{DSS}$ (la corriente con $V_{GS} = 0$), busque el $V_{GS}$ con el que
        $I_D = I_{DSS}/4$ y obtenga $V_P$ con~\eqref{eq:vp}.
  \item \label{paso:xy}Ponga \texttt{W1} como onda triangular de \qty{10}{\hertz} que vaya de
        \qty{0}{\volt} a $V_P$ (por ejemplo, amplitud \qty{1}{\volt} y offset
        \qty{-1}{\volt} si su $V_P$ es de unos \qty{-2}{\volt}). Pase el osciloscopio a modo
        \textbf{XY} con CH1 en el eje X y CH2 en el eje Y, y capture la curva.
\end{enumerate}

\noindent\begin{minipage}{\linewidth}
\paso{4}{Anote}
  \centering
  \renewcommand{\arraystretch}{1.15}
  \begin{tabularx}{\linewidth}{|>{\centering\arraybackslash}m{2.2cm}|*{4}{>{\centering\arraybackslash}X|}}
    \hline
    $\boldsymbol{V_{GS}}$ & \textbf{$I_D$ teórico} & \textbf{$I_D$ simulado}
      & \textbf{$v_{R_D}$ medido} & \textbf{$I_D$ medido} \\ \hline
    \qty{0}{\volt}     & \hueco & & & \\ \hline
    \qty{-0.25}{\volt} & \hueco & & & \\ \hline
    \qty{-0.5}{\volt}  & \hueco & & & \\ \hline
    \qty{-0.75}{\volt} & \hueco & & & \\ \hline
    \qty{-1.0}{\volt}  & \hueco & & & \\ \hline
    \qty{-1.25}{\volt} & \hueco & & & \\ \hline
    \hueco & \multicolumn{2}{c|}{} & & \\ \hline
    \hueco & \multicolumn{2}{c|}{} & & \\ \hline
  \end{tabularx}
  \captionof{table}{Experimento 1: curva de transferencia del JFET. Use las dos últimas filas
                    si su JFET necesita un $V_{GS}$ más negativo para cortarse.}
  \label{tab:exp1}
\end{minipage}\par\vspace{10pt}

\begin{table}[H]
  \centering
  \renewcommand{\arraystretch}{1.15}
  \begin{tabularx}{0.85\linewidth}{|>{\centering\arraybackslash}m{3.4cm}|*{3}{>{\centering\arraybackslash}X|}}
    \hline
    & \textbf{Hoja de datos} & \textbf{Simulado} & \textbf{Su JFET} \\ \hline
    $I_{DSS}$ & \qtyrange{1}{5}{\milli\ampere} & & \hueco \\ \hline
    $V_P$ con~\eqref{eq:vp} & \qtyrange{-0.5}{-6}{\volt} & & \hueco \\ \hline
  \end{tabularx}
  \caption{Experimento 1: los parámetros de su JFET. Los usará en el experimento~2.}
  \label{tab:param}
\end{table}

\noindent\begin{minipage}{\linewidth}
\noindent\textbf{Captura de la curva en modo XY} (paso~3\ref{paso:xy}).
\cuadrorespuesta
\end{minipage}\par

%-----------------------------------------------------------------------
\subsection{Experimento 2: el JFET como fuente de corriente para un LED}
%-----------------------------------------------------------------------
Un LED debe recibir siempre la misma corriente aunque el voltaje de alimentación cambie de
\qty{3}{\volt} a \qty{12}{\volt}. Se comparan tres circuitos (figura~\ref{fig:exp2}): el JFET
con la compuerta unida a la fuente (A), el JFET autopolarizado (B) y el LED con solo una
resistencia (C).

\paso{1}{Datos}
$V_{DD}$ de \qty{3}{}, \qty{4}{}, \qty{5}{} y \qty{12}{\volt}; LED rojo de \qty{1.85}{\volt} a
\qty{2}{\milli\ampere}. En el circuito B se desea $I_D = I_{DSS}/2$, con
$R_G = \qty{1}{\mega\ohm}$; en el circuito C, $R = \qty{4.7}{\kilo\ohm}$. En casa use
$I_{DSS} = \qty{2.1}{\milli\ampere}$ y $V_P = \qty{-1.4}{\volt}$; en el laboratorio, los de su
JFET (tabla~\ref{tab:param}).

\begin{figure}[H]
  \centering
  \begin{subfigure}[b]{0.30\linewidth}
    \centering
    \begin{circuitikz}[scale=0.8, transform shape]
      \draw (0,0) node[njfet](J){};
      \coordinate (top) at ($(J.D)+(0,2.4)$);
      \draw (top) to[leDo, l=LED] (J.D);
      \draw (top) node[circ]{} node[above]{$V_{DD}$};
      \draw (J.G) -- ++(-0.5,0) coordinate(gg) -- (gg |- 0,-1.4) -- (J.S |- 0,-1.4);
      \draw (J.S) -- (J.S |- 0,-1.4) node[circ]{} -- ++(0,-0.4) node[ground]{};
      \draw[->] ($(J.D)+(-0.35,1.0)$) -- node[left]{$I_{DSS}$} ++(0,-0.6);
    \end{circuitikz}
    \caption{A: $R_S = 0$.}
  \end{subfigure}
  \hfill
  \begin{subfigure}[b]{0.34\linewidth}
    \centering
    \begin{circuitikz}[scale=0.8, transform shape]
      \draw (0,0) node[njfet](J){};
      \coordinate (top) at ($(J.D)+(0,2.4)$);
      \draw (top) to[leDo, l=LED] (J.D);
      \draw (top) node[circ]{} node[above]{$V_{DD}$};
      \coordinate (gl) at (0,-2.6);
      \draw (J.S) to[R, l=$R_S$, a=?] (J.S |- gl) node[ground]{};
      \draw (J.G) -- ++(-0.6,0) coordinate(gg) to[R, l_=$R_G$] (gg |- gl) node[ground]{};
      \draw[->] ($(J.D)+(-0.35,1.0)$) -- node[left]{$I_D$} ++(0,-0.6);
    \end{circuitikz}
    \caption{B: autopolarizado.}
  \end{subfigure}
  \hfill
  \begin{subfigure}[b]{0.30\linewidth}
    \centering
    \begin{circuitikz}[scale=0.8, transform shape]
      \coordinate (top) at (0,2.4);
      \draw (top) to[leDo, l=LED] (0,0.6) to[R, l=$R$] (0,-1.6) node[ground]{};
      \draw (top) node[circ]{} node[above]{$V_{DD}$};
      \draw[->] (-0.35,1.9) -- node[left]{$I$} ++(0,-0.6);
    \end{circuitikz}
    \caption{C: solo resistencia.}
  \end{subfigure}
  \caption{Experimento 2: tres maneras de alimentar un LED. En A y B el JFET fija la corriente;
           en C la corriente depende de $V_{DD}$.}
  \label{fig:exp2}
\end{figure}

\paso{2}{Calcule (en casa)}
\begin{enumerate}[label=\alph*), leftmargin=*, itemsep=3pt]
  \item $V_{GS}$ y $R_S$ del circuito B con~\eqref{eq:rs}, como en el ejemplo de la
        tabla~\ref{tab:ejemplo}, y el valor E12 más cercano.
  \item $V_{DD(\min)}$ de los circuitos A y B con~\eqref{eq:vddmin}.
  \item La corriente del circuito C para cada $V_{DD}$.
  \item Simule los tres circuitos con un barrido DC de $V_{DD}$ de \qty{1}{\volt} a
        \qty{12}{\volt}.
\end{enumerate}

\paso{3}{Arme y mida (en el laboratorio)}
\begin{enumerate}[label=\alph*), ref=\alph*, leftmargin=*, itemsep=3pt]
  \item Recalcule $R_S$ con el $I_{DSS}$ y el $V_P$ que midió en su JFET
        (tabla~\ref{tab:param}) y elija el valor E12 más cercano de los que trajo.
  \item Arme el circuito A. Aliméntelo con la fuente programable $V_+$ en \qty{3}{},
        \qty{4}{} y \qty{5}{\volt} y luego con el riel de \qty{+12}{\volt}. En cada caso mida
        la corriente con el multímetro en serie con el LED y observe su brillo.
  \item Repita con el circuito B, donde $I_D = v_{R_S}/R_S$. Con \qty{12}{\volt} mida además
        $V_{GS}$, $V_{DS}$ y el voltaje en $R_G$.
  \item Repita con el circuito C.
\end{enumerate}

\noindent\begin{minipage}{\linewidth}
\paso{4}{Anote}
  \centering
  \renewcommand{\arraystretch}{1.15}
  \begin{tabularx}{\linewidth}{|>{\centering\arraybackslash}m{2.2cm}|>{\centering\arraybackslash}m{1.6cm}|*{4}{>{\centering\arraybackslash}X|}}
    \hline
    \multirow{2}{*}{\textbf{Circuito}} & & \multicolumn{4}{c|}{\textbf{Corriente en el LED}} \\
    \cline{3-6}
    & & $V_{DD} = \qty{3}{\volt}$ & \qty{4}{\volt} & \qty{5}{\volt} & \qty{12}{\volt} \\ \hline
    \multirow{2}{*}{A: $R_S = 0$}  & Simulado & \hueco & & & \\ \cline{2-6}
                                   & Medido   & \hueco & & & \\ \hline
    \multirow{2}{=}{\centering B: $R_S = $\\ \rule{1.4cm}{0.4pt}}
                                   & Simulado & \hueco & & & \\ \cline{2-6}
                                   & Medido   & \hueco & & & \\ \hline
    \multirow{2}{=}{\centering C: \qty{4.7}{\kilo\ohm}}
                                   & Teórico  & \hueco & & & \\ \cline{2-6}
                                   & Medido   & \hueco & & & \\ \hline
  \end{tabularx}
  \captionof{table}{Experimento 2: la corriente en el LED con cada voltaje de alimentación.}
  \label{tab:exp2}
\end{minipage}\par\vspace{10pt}

\begin{table}[H]
  \centering
  \renewcommand{\arraystretch}{1.15}
  \begin{tabularx}{0.85\linewidth}{|>{\centering\arraybackslash}m{4.2cm}|*{3}{>{\centering\arraybackslash}X|}}
    \hline
    \textbf{Circuito B con \qty{12}{\volt}} & \textbf{Teórico} & \textbf{Simulado} & \textbf{Medido} \\ \hline
    $V_{GS}$ & \hueco & & \\ \hline
    $I_D$    & \hueco & & \\ \hline
    $V_{DS}$ & \hueco & & \\ \hline
    $V_{R_G}$ (y entonces $I_G$) & \hueco & & \\ \hline
    $V_{DD(\min)}$ & \hueco & & \\ \hline
  \end{tabularx}
  \caption{Experimento 2: punto Q del JFET autopolarizado.}
  \label{tab:q2}
\end{table}

%-----------------------------------------------------------------------
\subsection{Experimento 3: el MOSFET de enriquecimiento}
%-----------------------------------------------------------------------
Comienza la semana~2. Se mide el umbral del 2N7000, su curva de transferencia y su resistencia
cuando conduce, y se comprueba que la compuerta no toma corriente.

\paso{1}{Datos}
\begin{table}[H]
  \centering
  \renewcommand{\arraystretch}{1.2}
  \begin{tabularx}{0.8\linewidth}{|>{\raggedright\arraybackslash}X|>{\centering\arraybackslash}m{5.4cm}|}
    \hline
    MOSFET & 2N7000 \\ \hline
    $V_{GS(\text{th})}$ y $K$ (en casa) & \qty{2.1}{\volt} y \qty{80}{\milli\ampere\per\volt\squared} \\ \hline
    $R_{DS(\text{on})}$ típica & \qty{2.5}{\ohm} con \qty{5}{\volt};
                                 \qty{1.2}{\ohm} con \qty{12}{\volt} \\ \hline
  \end{tabularx}
  \caption{Experimento 3: datos del 2N7000. Los valores de cada montaje están en la
           figura~\ref{fig:exp3}.}
  \label{tab:cond3}
\end{table}

\begin{figure}[H]
  \centering
  \begin{subfigure}[b]{0.42\linewidth}
    \centering
    \begin{circuitikz}[scale=0.85, transform shape]
      \draw (0,0) node[nigfete](M){};
      \node[right] at ($(M.S)+(0.05,0.35)$) {2N7000};
      \draw (M.D) -- ++(0,0.4) coordinate(dn) to[R, l_=$R_T$, a^=\qty{3.3}{\kilo\ohm}]
            ++(0,2.0) coordinate(top);
      \draw (top) node[circ]{} -- ++(1.2,0) node[right]{\qty{+5}{\volt}};
      \draw (M.G) -- ++(-0.7,0) coordinate(gg) -- (gg |- dn) -- (dn) node[circ]{};
      \draw (dn) to[short, -o] ++(1.2,0) node[right]{$V_{GS(\text{th})}$};
      \draw (M.S) node[ground]{};
    \end{circuitikz}
    \caption{Umbral: compuerta unida al drenador.}
  \end{subfigure}
  \hfill
  \begin{subfigure}[b]{0.54\linewidth}
    \centering
    \begin{circuitikz}[scale=0.85, transform shape]
      \draw (0,0) node[nigfete](M){};
      \node[right] at ($(M.S)+(0.05,0.35)$) {2N7000};
      \draw (M.D) -- ++(0,0.3) coordinate(d) to[R, l=$R_D$] ++(0,2.0) coordinate(top);
      \draw (top) node[circ]{} -- ++(1.2,0) node[right]{$V_{DD}$};
      \draw (d) to[short, *-o] ++(1.2,0) node[right]{$D$};
      \draw (M.G) to[R, l_=$R_G$, a^=\qty{100}{\kilo\ohm}] ++(-2.6,0)
            node[ocirc]{} node[left]{\texttt{W1}};
      \draw (M.S) node[ground]{};
    \end{circuitikz}
    \caption{Transferencia y $R_{DS(\text{on})}$.}
  \end{subfigure}
  \caption{Experimento 3: circuitos de medida del 2N7000. En (b), $V_{DD} = \qty{12}{\volt}$ y
           $R_D = \qty{1}{\kilo\ohm}$ para la curva; $V_{DD} = \qty{5}{\volt}$ y
           $R_D = \qty{100}{\ohm}$ (\qty{0.5}{\watt}) para $R_{DS(\text{on})}$.}
  \label{fig:exp3}
\end{figure}

\paso{2}{Calcule (en casa)}
\begin{enumerate}[label=\alph*), leftmargin=*, itemsep=3pt]
  \item La corriente del circuito de umbral si $V_{GS} = V_{GS(\text{th})}$. ¿Se parece a la
        \qty{1}{\milli\ampere} con la que la hoja de datos define el umbral?
  \item $I_D$ con~\eqref{eq:mos} para los $V_{GS}$ de la tabla~\ref{tab:exp3}. La corriente
        no puede pasar de $V_{DD}/R_D = \qty{12}{\milli\ampere}$: ¿con qué $V_{GS}$ se llega a
        ese límite? Desde ahí el MOSFET está en la región óhmica.
  \item $I_D$ y $V_{DS}$ con~\eqref{eq:rdson} en la medida de $R_{DS(\text{on})}$, con
        $V_{GS} = \qty{5}{\volt}$ y \qty{12}{\volt}.
  \item Simule los tres casos.
\end{enumerate}

\paso{3}{Arme y mida (en el laboratorio)}
\begin{enumerate}[label=\alph*), ref=\alph*, leftmargin=*, itemsep=3pt]
  \item Arme el circuito de umbral con el riel de \qty{+5}{\volt} y mida $V_{GS(\text{th})}$
        y el voltaje en $R_T$.
  \item Arme el circuito de la figura~\ref{fig:exp3}b con \qty{12}{\volt} y
        $R_D = \qty{1}{\kilo\ohm}$, con la misma conexión de la figura~\ref{fig:conexion}.
        Ponga \texttt{W1} en continua con cada $V_{GS}$ de la tabla~\ref{tab:exp3} y mida
        $v_{R_D}$ y $V_{DS}$.
  \item \label{paso:xymos}Ponga \texttt{W1} como onda triangular de \qty{10}{\hertz} de
        \qty{0}{\volt} a \qty{5}{\volt} (amplitud \qty{2.5}{\volt}, offset
        \qty{2.5}{\volt}) y capture la curva en modo XY.
  \item Con \texttt{W1} en \qty{5}{\volt}, mida el voltaje en $R_G$. ¿Cuánta corriente toma la
        compuerta?
  \item Cambie a $V_{DD} = \qty{5}{\volt}$ y $R_D = \qty{100}{\ohm}$. Con \texttt{W1} en
        \qty{5}{\volt} mida $V_{DS}$ y $v_{R_D}$. Luego pase el extremo libre de $R_G$ de
        \texttt{W1} al riel de \qty{+12}{\volt} ($V_{GS} = \qty{12}{\volt}$) y repita.
\end{enumerate}

\noindent\begin{minipage}{\linewidth}
\paso{4}{Anote}
  \centering
  \renewcommand{\arraystretch}{1.15}
  \begin{tabularx}{\linewidth}{|>{\centering\arraybackslash}m{2.0cm}|*{5}{>{\centering\arraybackslash}X|}}
    \hline
    $\boldsymbol{V_{GS}}$ & \textbf{$I_D$ teórico} & \textbf{$I_D$ simulado}
      & \textbf{$I_D$ medido} & \textbf{$V_{DS}$ medido} & \textbf{Región} \\ \hline
    \qty{1.5}{\volt} & \hueco & & & & \\ \hline
    \qty{2.0}{\volt} & \hueco & & & & \\ \hline
    \qty{2.2}{\volt} & \hueco & & & & \\ \hline
    \qty{2.4}{\volt} & \hueco & & & & \\ \hline
    \qty{2.6}{\volt} & \hueco & & & & \\ \hline
    \qty{3.0}{\volt} & \hueco & & & & \\ \hline
    \qty{4.0}{\volt} & \hueco & & & & \\ \hline
    \qty{5.0}{\volt} & \hueco & & & & \\ \hline
  \end{tabularx}
  \captionof{table}{Experimento 3: curva de transferencia del 2N7000 con
                    $V_{DD} = \qty{12}{\volt}$ y $R_D = \qty{1}{\kilo\ohm}$.}
  \label{tab:exp3}
\end{minipage}\par\vspace{10pt}

\begin{table}[H]
  \centering
  \renewcommand{\arraystretch}{1.15}
  \begin{tabularx}{\linewidth}{|>{\centering\arraybackslash}m{5.0cm}|*{3}{>{\centering\arraybackslash}X|}}
    \hline
    \textbf{Parámetro} & \textbf{Teórico} & \textbf{Simulado} & \textbf{Medido} \\ \hline
    $V_{GS(\text{th})}$ e $I_D$ en el umbral & \hueco & & \\ \hline
    $I_G$ con $V_{GS} = \qty{5}{\volt}$ & \hueco & & \\ \hline
    $V_{DS}$ e $I_D$ con $V_{GS} = \qty{5}{\volt}$ & \hueco & & \\ \hline
    $R_{DS(\text{on})}$ con $V_{GS} = \qty{5}{\volt}$ & \hueco & & \\ \hline
    $V_{DS}$ e $I_D$ con $V_{GS} = \qty{12}{\volt}$ & \hueco & & \\ \hline
    $R_{DS(\text{on})}$ con $V_{GS} = \qty{12}{\volt}$ & \hueco & & \\ \hline
  \end{tabularx}
  \caption{Experimento 3: umbral, corriente de compuerta y resistencia de conducción.}
  \label{tab:rdson}
\end{table}

\noindent\begin{minipage}{\linewidth}
\noindent\textbf{Captura de la curva del MOSFET en modo XY} (paso~3\ref{paso:xymos}).
\cuadrorespuesta
\end{minipage}\par

%-----------------------------------------------------------------------
\subsection{Experimento 4: atenuador PWM de LEDs}
%-----------------------------------------------------------------------
Aplicación del MOSFET como interruptor: una señal PWM de \qty{5}{\volt}, como la de un
microcontrolador o de un temporizador 555, regula el brillo de seis LED alimentados con
\qty{12}{\volt}. Aquí la señal la da \texttt{W1}.

\noindent\begin{minipage}{\linewidth}
\paso{1}{Datos}
\noindent\begin{minipage}[c]{0.44\linewidth}
  \centering
  \renewcommand{\arraystretch}{1.1}
  \begin{tabularx}{\linewidth}{|>{\raggedright\arraybackslash}X|>{\centering\arraybackslash}m{2.5cm}|}
    \hline
    $V_{DD}$ & \qty{12}{\volt} \\ \hline
    Cada rama & 3 LED rojos y \qty{330}{\ohm} \\ \hline
    $V_{\text{LED}}$ a \qty{18}{\milli\ampere} & \qty{2.0}{\volt} \\ \hline
    PWM (\texttt{W1}) & \qty{0}{} a \qty{5}{\volt}, \qty{1}{\kilo\hertz} \\ \hline
    $R_G$ y $R_{GS}$ & \qty{1}{\kilo\ohm} y \qty{100}{\kilo\ohm} \\ \hline
  \end{tabularx}
  \captionof{table}{Experimento 4: datos. Use la $R_{DS(\text{on})}$ que midió en el
                    experimento~3.}
  \label{tab:cond4}
\end{minipage}\hfill
\begin{minipage}[c]{0.53\linewidth}
  \centering
  \begin{circuitikz}[scale=0.66, transform shape]
    \foreach \x in {0,2.2}{
      \draw (\x,5.6) to[leDo] (\x,4.6) to[leDo] (\x,3.6) to[leDo] (\x,2.6)
            to[R, l=\qty{330}{\ohm}] (\x,1.2);
    }
    \draw (0,5.6) -- (2.2,5.6);
    \draw (1.1,5.6) node[circ]{} -- ++(0,0.4) node[above]{$V_{DD}$};
    \draw (0,1.2) -- (2.2,1.2);
    \draw (1.1,0.6) node[nigfete, anchor=D](M){};
    \draw (1.1,1.2) node[circ]{} -- (M.D);
    \node[right] at ($(M.S)+(0.05,0.35)$) {2N7000};
    \draw (2.2,1.2) to[short, -o] ++(1.0,0) node[right]{$D$};
    \draw (M.G) -- ++(-0.5,0) coordinate(gn) to[R, l_=$R_G$]
          ++(-2.0,0) node[ocirc]{} node[left]{\texttt{W1}};
    \draw (gn) node[circ]{} to[R, l=$R_{GS}$] ($(gn |- M.S)+(0,-1.0)$)
          -- ($(M.S)+(0,-1.0)$);
    \draw (M.S) -- ++(0,-1.0) node[circ]{} node[ground]{};
  \end{circuitikz}
  \captionof{figure}{Experimento 4: atenuador PWM.}
  \label{fig:exp4}
\end{minipage}
\end{minipage}\par\vspace{6pt}

\paso{2}{Calcule (en casa)}
\begin{enumerate}[label=\alph*), leftmargin=*, itemsep=3pt]
  \item $I_{on}$ de cada rama y la total con el MOSFET conduciendo ($V_{DS} \approx 0$), la
        potencia en cada resistencia de \qty{330}{\ohm} y la potencia en el MOSFET con
        $I_{on}^2 R_{DS(\text{on})}$.
  \item $I_{\text{prom}}$ total con~\eqref{eq:pwm} para cada ciclo de trabajo de la
        tabla~\ref{tab:exp4}.
  \item Simule el circuito con una fuente de pulsos de \qty{1}{\kilo\hertz} y mida la
        corriente promedio para cada ciclo de trabajo.
\end{enumerate}

\paso{3}{Arme y mida (en el laboratorio)}
\begin{enumerate}[label=\alph*), ref=\alph*, leftmargin=*, itemsep=3pt]
  \item Arme el circuito con el riel de \qty{+12}{\volt}. $R_{GS}$ mantiene apagado el MOSFET
        si \texttt{W1} se desconecta y $R_G$ protege la salida de \texttt{W1}: no las omita.
        Ponga \texttt{W1} en \qty{5}{\volt} continuos y mida $V_{DS}$ y la corriente total con
        el multímetro en serie con el riel.
  \item Ponga \texttt{W1} como onda cuadrada de \qty{1}{\kilo\hertz}, amplitud
        \qty{2.5}{\volt} y offset \qty{2.5}{\volt}. Para cada ciclo de trabajo (\emph{Symmetry}
        en WaveForms) de la tabla~\ref{tab:exp4}, mida la corriente promedio con el multímetro
        en DC y califique el brillo de 1 (casi apagado) a 5 (máximo).
  \item \label{paso:pwm}Con $D = \qty{25}{\percent}$, capture $v_{GS}$ (CH1) y $v_{DS}$ (CH2).
  \item Con $D = \qty{50}{\percent}$, baje la frecuencia a \qty{5}{}, \qty{20}{}, \qty{50}{}
        y \qty{100}{\hertz} y anote desde qué frecuencia deja de verse el parpadeo.
\end{enumerate}

\noindent\begin{minipage}{\linewidth}
\paso{4}{Anote}
  \centering
  \renewcommand{\arraystretch}{1.15}
  \begin{tabularx}{\linewidth}{|>{\centering\arraybackslash}m{2.4cm}|*{4}{>{\centering\arraybackslash}X|}}
    \hline
    \textbf{Ciclo de trabajo} & \textbf{$I_{\text{prom}}$ teórica} & \textbf{$I_{\text{prom}}$ simulada}
      & \textbf{$I_{\text{prom}}$ medida} & \textbf{Brillo (1 a 5)} \\ \hline
    \qty{10}{\percent} & \hueco & & & \\ \hline
    \qty{25}{\percent} & \hueco & & & \\ \hline
    \qty{50}{\percent} & \hueco & & & \\ \hline
    \qty{75}{\percent} & \hueco & & & \\ \hline
    \qty{90}{\percent} & \hueco & & & \\ \hline
    \qty{100}{\percent} (\texttt{W1} en \qty{5}{\volt}) & \hueco & & & \\ \hline
  \end{tabularx}
  \captionof{table}{Experimento 4: corriente promedio y brillo contra el ciclo de trabajo.}
  \label{tab:exp4}
\end{minipage}\par\vspace{10pt}

\begin{table}[H]
  \centering
  \renewcommand{\arraystretch}{1.15}
  \begin{tabularx}{\linewidth}{|>{\centering\arraybackslash}m{3.4cm}|*{4}{>{\centering\arraybackslash}X|}}
    \hline
    \textbf{Frecuencia} & \qty{5}{\hertz} & \qty{20}{\hertz} & \qty{50}{\hertz} & \qty{100}{\hertz} \\ \hline
    ¿Se ve el parpadeo? & \hueco & & & \\ \hline
  \end{tabularx}
  \caption{Experimento 4: el parpadeo con $D = \qty{50}{\percent}$.}
  \label{tab:parpadeo}
\end{table}

\noindent\begin{minipage}{\linewidth}
\noindent\textbf{Captura de $v_{GS}$ y $v_{DS}$ con $D = \qty{25}{\percent}$} (paso~3\ref{paso:pwm}).
\cuadrorespuesta
\end{minipage}\par

\FloatBarrier
%=======================================================================
\section{Informe final}
%=======================================================================
\begin{enumerate}[label=\textbf{\arabic*.}, leftmargin=*, labelsep=0.6em]
  \item Adjunte los cálculos y las simulaciones de los cuatro experimentos.
  \item Compare los valores teóricos, simulados y medidos y explique las diferencias más
        grandes. Grafique en un mismo eje la curva de transferencia medida de su JFET y la de
        la ecuación de Shockley con su $I_{DSS}$ y su $V_P$.
  \item Explique con la tabla~\ref{tab:exp2} por qué el JFET mantiene constante la corriente
        del LED y la resistencia sola no.
  \item Compare con la tabla~\ref{tab:bjtmos} el interruptor con MOSFET del experimento~4 y el
        interruptor con BJT del Laboratorio~6.
\end{enumerate}

\FloatBarrier
%=======================================================================
\section{Cuestionario}
%=======================================================================
\begin{enumerate}[label=\textbf{\arabic*.}, leftmargin=*, labelsep=0.6em]
\item En el experimento 2, demuestre con~\eqref{eq:autopol} y~\eqref{eq:vddmin} que el voltaje
mínimo de alimentación de la fuente de corriente es $V_{\text{LED}} + |V_P|$, sin importar el
valor de $R_S$. ¿Cuánto vale con el $V_P$ de su JFET? ¿Lo confirma la tabla~\ref{tab:exp2}?
\lineasresp{3}

\item En el Laboratorio~6 el relé (\qty{68.6}{\milli\ampere}) necesitó
$I_B = \qty{3.6}{\milli\ampere}$ y el 2N2222A cayó $V_{CE(\text{sat})} = \qty{0.12}{\volt}$.
Si lo maneja un 2N7000 con $V_{GS} = \qty{5}{\volt}$ y la $R_{DS(\text{on})}$ que usted midió,
¿cuánta corriente toma la compuerta en DC? ¿Cuánto valen $V_{DS}$ y la potencia en el
transistor en cada caso?
\lineasresp{3}

\item En el experimento 4, con $D = \qty{30}{\percent}$, ¿qué corriente promedio y qué potencia
promedio entrega el riel de \qty{12}{\volt}? ¿Por qué a \qty{5}{\hertz} se ve el parpadeo y a
\qty{1}{\kilo\hertz} no? Según su tabla~\ref{tab:exp4}, ¿el brillo que se percibe es
proporcional a $D$?
\lineasresp{3}
\end{enumerate}

\FloatBarrier
%=======================================================================
\section{Conclusiones}
%=======================================================================
\begin{itemize}
    \item Redacte como mínimo 05 conclusiones \textbf{cualitativas} y 05
          \textbf{cuantitativas}.
\end{itemize}

\noindent\textbf{Cualitativas}
\lineasresp{10}

\needspace{6\baselineskip}
\noindent\textbf{Cuantitativas}
\lineasresp{10}

\end{document}
